% Continues Remark 3.5.1.
One will denote by \(\delta(\mu,\mu')\) the image of \(\delta(\mu,\mu')\) in
\(H^{2}(X_{0},L_{0})\). Finally, if \(X\) equipped with the group law \(\mu\)
(\resp \(\mu'\)) is denoted simply by \(X\) (\resp \(X'\)), one will write
\(\delta(X,X')\) instead of \(\delta(\mu,\mu')\), and likewise for
\(\delta(X,X')\).

\begin{remark}{3.6}\label{III.3.6}
Let \(X_{J}\) be an \(S_{J}\)-scheme \emph{smooth} over \(S_{J}\) and
\emph{affine}. By~0.15, there exists, up to isomorphism, a unique
\(S\)-scheme \(X\), \emph{smooth} over \(S\), and reducing to \(X_{J}\). If
\(X_{J}\) is equipped with a structure of \(S_{J}\)-group, it follows
from~0.16 that condition \((i_{1})\) is automatically satisfied. Moreover,
by~0.6 the definition of \(L_{0}\) simplifies and one obtains:
\end{remark}

\begin{corollary}{3.7}\label{III.3.7}
Let \(S\), \(\mathcal{I}\) and \(\mathcal{J}\) be as in~3.1. Let \(X_{J}\)
be an \(S_{J}\)-group smooth over \(S_{J}\) and affine.
\begin{enumeratei}
\item The set of \(S\)-groups smooth over \(S\) and reducing to \(X_{J}\),
up to isomorphism (inducing the identity on \(X_{J}\)), is empty or a
principal homogeneous set under the group
\[
H^{2}\bigl(X_{0},\mathcal{L}\mathrm{ie}(X_{0}/S_{0})
  \otimes_{\mathcal{O}_{S_{0}}}\mathcal{J}\bigr).
\]
\item There exists an \(S\)-group smooth over \(S\) reducing to \(X_{J}\) if
and only if a certain obstruction in
\[
H^{3}\bigl(X_{0},\mathcal{L}\mathrm{ie}(X_{0}/S_{0})
  \otimes_{\mathcal{O}_{S_{0}}}\mathcal{J}\bigr)
\]
is zero.
\end{enumeratei}
\end{corollary}

One deduces from this, as usual, the following corollaries:

\begin{corollary}{3.8}\label{III.3.8}
Let \(S\) be a scheme and \(S_{0}\) a closed subscheme defined by a
nilpotent ideal \(\mathcal{I}\). Let \(X_{0}\) be an \(S_{0}\)-group smooth
over \(S\) and affine.%
% typo?: kept as printed (lisse sur S, not sur S_0)
\begin{enumeratei}
\item If
\(H^{2}\bigl(X_{0},\mathcal{L}\mathrm{ie}(X_{0}/S_{0})
  \otimes_{\mathcal{O}_{S_{0}}}\mathcal{I}^{n+1}/\mathcal{I}^{n+2}\bigr)=0\)
for every \(n\geqslant 0\), two \(S\)-groups smooth over \(S\) and reducing
to \(X_{0}\) are isomorphic (by an isomorphism inducing the identity on
\(X_{0}\)).
\item If
\(H^{3}\bigl(X_{0},\mathcal{L}\mathrm{ie}(X_{0}/S_{0})
  \otimes_{\mathcal{O}_{S_{0}}}\mathcal{I}^{n+1}/\mathcal{I}^{n+2}\bigr)=0\)
for every \(n\geqslant 0\), there exists an \(S\)-group smooth over \(S\),
reducing to \(X_{0}\).
\end{enumeratei}
\end{corollary}

\begin{corollary}{3.9}\label{III.3.9}
Let \(S\) be an affine scheme and \(S_{0}\) a closed subscheme defined by a
nilpotent ideal \(\mathcal{I}\). Suppose the
\(\mathcal{I}^{n+1}/\mathcal{I}^{n+2}\) are locally free over \(S_{0}\).
Let \(X_{0}\) be an \(S_{0}\)-group smooth and affine over
\(S_{0}\).%
% typo: missing period after S_0
\begin{enumeratei}
\item If \(H^{2}\bigl(X_{0},\mathcal{L}\mathrm{ie}(X_{0}/S_{0})\bigr)=0\),
two \(S\)-groups smooth over \(S\) and reducing to \(X_{0}\) are isomorphic.
\item If \(H^{3}\bigl(X_{0},\mathcal{L}\mathrm{ie}(X_{0}/S_{0})\bigr)=0\),
there exists an \(S\)-group smooth over \(S\) reducing to \(X_{0}\).
\end{enumeratei}
\end{corollary}

\begin{corollary}{3.10}\label{III.3.10}
Let \(S_{0}\) be a scheme and \(S=\mathrm{I}_{S_{0}}\) the scheme of dual
numbers over \(S_{0}\). Let \(X_{0}\) be an \(S_{0}\)-group smooth over
\(S_{0}\). In order that every \(S\)-group \(Y\), smooth over \(S\), such
that \(Y_{0}\) is \(S_{0}\)-isomorphic to \(X_{0}\), be \(S\)-isomorphic to
\(X=X_{0}\times_{S_{0}}S\), it is necessary and sufficient that
\(H^{2}\bigl(X_{0},\mathcal{L}\mathrm{ie}(X_{0}/S_{0})\bigr)=0\).%
{\renewcommand{\thefootnote}{(72)}\nde{This is used in XXIV, 1.13.}}
\end{corollary}

Indeed, by virtue of~3.5, the set of classes, up to an isomorphism of
\(S\)-groups ``inducing the identity on \(X_{0}\)'', of such groups \(Y\),
is in bijection with
\(H^{2}\bigl(X_{0},\mathcal{L}\mathrm{ie}(X_{0}/S_{0})\bigr)\), hence the
set of classes, up to an arbitrary isomorphism of \(S\)-groups, is in
bijection with
\[
H^{2}\bigl(X_{0},\mathcal{L}\mathrm{ie}(X_{0}/S_{0})\bigr)/\Gamma_{0}\,,
\]
where
\[
\Gamma_{0}=\Aut_{S_{0}\text{-gr.}}(X_{0})
\]
(which operates in an evident way on the \(H^{2}\)). The conclusion follows
at once from this.%
{\renewcommand{\thefootnote}{(73)}\nde{Indeed,
\(\Aut_{S_{0}\text{-gr.}}(X_{0})\) operates by group automorphisms on the
abelian group
\(H^{2}\bigl(X_{0},\mathcal{L}\mathrm{ie}(X_{0}/S_{0})\bigr)\), hence the
orbit of \(0\) is the singleton \(\{0\}\); consequently the quotient set is
a singleton if and only if
\(H^{2}\bigl(X_{0},\mathcal{L}\mathrm{ie}(X_{0}/S_{0})\bigr)=\{0\}\).}}

\sgasection{4}{Infinitesimal extensions of closed subgroups}
\label{III.sec.4}

Let us first state a result valid in an arbitrary \emph{abelian category}.

\begin{lemma}{4.1}\label{III.4.1}
Let \(0\to A'\xrightarrow{i}A\xrightarrow{p}A''\to 0\) be an exact sequence,
\(\phi\colon A'\to Q\) a morphism and \(\pi\colon A''\to P\) an epimorphism
with kernel \(C\). Let \(E\) be the set (up to isomorphism) of the
quadruples \((B,f,g,h)\) such that the sequence
\[
0 \longrightarrow Q \xrightarrow{f} B \xrightarrow{g} P \longrightarrow 0
\]
is exact and the diagram below is commutative:
\[
\xymatrix@C=2.6em@R=2.4em{
0 \ar[r] & A' \ar[r]^{i} \ar[d]_{\phi}
  & A \ar[r]^{p} \ar[d]_{h}
  & A'' \ar[r] \ar[d]_{\pi} & 0 \\
0 \ar[r] & Q \ar[r]^{f}
  & B \ar[r]^{g}
  & P \ar[r] & 0.
}
\]
\begin{enumeratei}
\item In order that \(E\) be non-empty, it is necessary and sufficient that
the image in \(\Ext^{1}(C,Q)\) of the element \(A\) of \(\Ext^{1}(A'',A')\)
be zero.
\item Under these conditions, \(E\) is a principal homogeneous set under
the abelian group \(\Hom(C,Q)\).
\end{enumeratei}
\end{lemma}

Let us introduce the amalgamated sum \(B'=A\amalg^{A'}Q\). One then has a
commutative diagram in which the rows are exact:%
{\renewcommand{\thefootnote}{(74)}\nde{One has
\(\Coker(j)=B'\amalg^{Q}0=A\amalg^{A'}0=A''\), and one sees that
\(\Ker(j)\simeq\Ker(i)=0\) by arguing ``as if \(C\) were a category of
modules''; for a proof solely in terms of arrows, see for example
\textup{[Fr64]}, Th.~2.5.4~$(*)$.}}
\[
\xymatrix@C=2.6em@R=2.4em{
0 \ar[r] & A' \ar[r]^{i} \ar[d]_{\phi}
  & A \ar[r]^{p} \ar[d]
  & A'' \ar[r] \ar[d] & 0 \\
0 \ar[r] & Q \ar[r]^{j}
  & B' \ar[r]
  & A'' \ar[r] & 0,
}
\]
and it is clear that the category of solutions of the problem posed is
canonically isomorphic to the category of solutions of the corresponding
problem for the sequence
\[
0 \longrightarrow Q \longrightarrow B' \longrightarrow A'' \longrightarrow 0
\]
and the morphisms \(\mathrm{id}_{Q}\) and \(\pi\colon A''\to P\).%
{\renewcommand{\thefootnote}{(75)}\nde{In what follows, one has replaced
\(A\) by \(B'\), and set out the end of the argument.}}
In this case, the set \(E\) is in bijection with the set of subobjects \(N\)
of \(B'\) such that \(B'\to A''\) induces an isomorphism of \(N\) with the
kernel \(C\) of \(A''\to P\), that is, the set of morphisms
\(\varepsilon\colon C\to B'\) lifting the canonical morphism \(C\to A''\).
The abelian group \(G=\Hom(C,Q)\) acts on \(E\) by
\(g\cdot\varepsilon=g+\varepsilon\) (addition in \(\Hom(C,B')\)), and if
\(E\neq\varnothing\) this makes \(E\) into a principal homogeneous set under
\(G\).

One deduces from this:

\begin{proposition}{4.2}\label{III.4.2}
{\renewcommand{\thefootnote}{(76)}\nde{One has rewritten the statement so
as to lie exactly in the setting of the application which is made of it
in~4.3; on the other hand, one has set out the proof, following the
indications given by M.~Demazure.}}%
Let \(S\) be a scheme, \(S_{J}\) the closed subscheme defined by a
quasi-coherent ideal \(\mathcal{J}\) of square zero, \(X\) an \(S\)-scheme,
\(\mathcal{F}\) an \(\mathcal{O}_{X}\)-module,
\(X_{J}=X\times_{S}S_{J}\),
\(\mathcal{F}_{J}=\mathcal{F}\otimes_{\mathcal{O}_{S}}\mathcal{O}_{S_{J}}\),
and \(\mathcal{G}_{J}=\mathcal{F}_{J}/\mathcal{H}_{J}\) a quotient module of
\(\mathcal{F}_{J}\). Suppose given a morphism of
\(\mathcal{O}_{X_{J}}\)-modules
\[
f\colon\mathcal{J}\otimes_{\mathcal{O}_{S_{J}}}\mathcal{G}_{J}
  \longrightarrow\mathcal{Q}.
\]
Let \(\mathcal{E}\) be the sheaf of sets on \(X\) defined as follows: for
each open \(U\) of \(X\), \(\mathcal{E}(U)\) is the set of quotient modules
\(\mathcal{G}\) of \(\mathcal{F}|_{U}\), such that
\(\mathcal{G}/\mathcal{J}\mathcal{G}=\mathcal{G}_{J}|_{U}\) and such that
there exists an isomorphism
\[
h\colon\mathcal{J}\mathcal{G}\isomto\mathcal{Q}|_{U}
\]
making the following diagram commutative:
\[
\xymatrix@C=1.2em@R=1.6em{
\mathcal{J}\otimes_{\mathcal{O}_{S_{J}}}(\mathcal{G}_{J}|_{U})
  \ar[rr]^{f|_{U}} \ar[dr]_{\mathrm{can.}}
&& \mathcal{Q}|_{U} \\
& \mathcal{J}\mathcal{G} \ar[ur]_{h}^{\sim}
}\,.
\]
(\(h\) is then unique, since
\(\mathcal{J}\otimes_{\mathcal{O}_{S_{J}}}(\mathcal{G}_{J}|_{U})
  \to\mathcal{J}\mathcal{G}\)
is an epimorphism). Then \(\mathcal{E}\) is a formally principal homogeneous
sheaf under the sheaf of commutative groups
\[
\mathcal{A}
  =\mathcal{H}\mathrm{om}_{\mathcal{O}_{X}}(\mathcal{H}_{J},\mathcal{Q})
  =\mathcal{H}\mathrm{om}_{\mathcal{O}_{X_{J}}}(\mathcal{H}_{J},\mathcal{Q}).
\]
\end{proposition}

\begin{proof}
If \(\mathcal{E}(U)=\varnothing\) there is nothing to prove; one may
therefore suppose that \(\mathcal{E}(U)\) contains an element
\(\widetilde{\mathcal{G}}\). Then, in the diagram below, \(h\) is an
isomorphism and all the arrows are epimorphisms:
\[
\xymatrix@C=1.6em@R=2.2em{
\mathcal{J}\otimes_{\mathcal{O}_{S_{J}}}(\mathcal{F}_{J}|_{U})
  \ar[r] \ar[d]_{\mathrm{can.}}
& \mathcal{J}\otimes_{\mathcal{O}_{S_{J}}}(\mathcal{G}_{J}|_{U})
  \ar[r]^{f|_{U}} \ar[d]_{\mathrm{can.}}
& \mathcal{Q}|_{U} \\
\mathcal{J}\mathcal{F}|_{U} \ar[r]
& \mathcal{J}\widetilde{\mathcal{G}} \ar[ur]_{h}^{\sim}
&
}\,.
\]
Thus, the morphism
\(\mathcal{J}\otimes_{\mathcal{O}_{S_{J}}}(\mathcal{F}_{J}|_{U})
  \to\mathcal{Q}|_{U}\)
induces an epimorphism (necessarily unique)
\(\phi\colon\mathcal{J}\mathcal{F}|_{U}\to\mathcal{Q}|_{U}\), and if
\(\mathcal{G}\) is an \(\mathcal{O}_{U}\)-module such that
\(\mathcal{G}/\mathcal{J}\mathcal{G}=\mathcal{G}_{J}|_{U}\) and one has a
commutative diagram with exact rows:
\[
\xymatrix@C=2.2em@R=2.2em{
0 \ar[r]
  & \mathcal{J}\mathcal{F}|_{U} \ar[r] \ar[d]_{\phi}
  & \mathcal{F}|_{U} \ar[r] \ar[d]
  & \mathcal{F}_{J}|_{U} \ar[r] \ar[d]_{\pi}
  & 0 \\
0 \ar[r]
  & \mathcal{Q}|_{U} \ar[r]
  & \mathcal{G} \ar[r]^{p_{J}}
  & \mathcal{G}_{J}|_{U} \ar[r]
  & 0
}
\]
(where \(p_{J}\) is the projection
\(\mathcal{G}\to\mathcal{G}/\mathcal{J}\mathcal{G}=\mathcal{G}_{J}|_{U}\),
so that \(\mathcal{Q}|_{U}=\Ker(p_{J})=\mathcal{J}\mathcal{G}\)), then one
may identify \(\mathcal{G}\) with a quotient module of
\(\mathcal{F}|_{U}\). Consequently, by~4.1~(ii), the set
\(\mathcal{E}(U)\) is a principal homogeneous set under the abelian group
\[
\mathcal{H}\mathrm{om}_{\mathcal{O}_{X}}(\mathcal{H}_{J},\mathcal{Q})(U)
  =\mathcal{H}\mathrm{om}_{\mathcal{O}_{X_{J}}}(\mathcal{H}_{J},\mathcal{Q})(U).
\]
\end{proof}

\begin{proposition}{4.3}\label{III.4.3}
(TDTE~IV~5.1) Let \(S\) be a scheme, \(S_{J}\) the closed subscheme defined
by a quasi-coherent ideal \(\mathcal{J}\) of square zero, \(X\) an
\(S\)-scheme, \(\mathcal{F}\) a quasi-coherent \(\mathcal{O}_{X}\)-module,
\(X_{J}=X\times_{S}S_{J}\),
\(\mathcal{F}_{J}=\mathcal{F}\otimes_{\mathcal{O}_{S}}\mathcal{O}_{S_{J}}\).
Let \(\mathcal{G}_{J}=\mathcal{F}_{J}/\mathcal{H}_{J}\) be a quasi-coherent
quotient module of \(\mathcal{F}_{J}\), flat over \(S_{J}\).

For every open \(U\) of \(X\), let \(\mathcal{E}(U)\) be the set of
quasi-coherent quotient modules \(\mathcal{G}\)%
{\renewcommand{\thefootnote}{(77)}\nde{In order to lighten the statement,
one has added here the hypothesis that \(\mathcal{G}\) be quasi-coherent,
and deferred to the proof the remark that this hypothesis is automatically
verified; one has set out the proof accordingly.}}
of \(\mathcal{F}|_{U}\), flat over \(S\) and such that
\(\mathcal{G}/\mathcal{J}\mathcal{G}\simeq\mathcal{G}_{J}|_{U}\). Then the
\(\mathcal{E}(U)\) form a sheaf of sets \(\mathcal{E}\) on \(X\), which is
formally principal homogeneous under the sheaf of commutative groups
\[
\mathcal{A}
  =\mathcal{H}\mathrm{om}_{\mathcal{O}_{X_{J}}}\bigl(\mathcal{H}_{J},
    \mathcal{J}\otimes_{\mathcal{O}_{S_{J}}}\mathcal{G}_{J}\bigr).
\]
\end{proposition}

\begin{proof}
Let us denote by \(\pi\colon X\to S\) the structural morphism. Let \(U\) be
an open of \(X\) and \(\mathcal{G}\) an \(\mathcal{O}_{U}\)-module
\emph{flat} over \(S\) and such that
\(\mathcal{G}/\mathcal{J}\mathcal{G}\simeq\mathcal{G}_{J}|_{U}\). Then, for
every \(x\in U\), \(\mathcal{G}_{x}\) is a module \emph{flat} over the local
ring \(\mathcal{O}_{S,s}\) (where \(s=\pi(x)\)), and therefore the morphism
\[
\mathcal{J}_{s}\otimes_{\mathcal{O}_{S,s}}(\mathcal{G}/\mathcal{J}\mathcal{G})_{x}
  =\mathcal{J}_{s}\otimes_{\mathcal{O}_{S,s}}\mathcal{G}_{x}
  \longrightarrow(\mathcal{J}\mathcal{G})_{x}
\]
is bijective; one therefore has an exact sequence
\[
0 \longrightarrow
  \mathcal{J}\otimes_{\mathcal{O}_{S}}(\mathcal{G}_{J}|_{U})
  \longrightarrow\mathcal{G}
  \longrightarrow\mathcal{G}_{J}|_{U}
  \longrightarrow 0
\]
and since \(\mathcal{J}\otimes_{\mathcal{O}_{S}}(\mathcal{G}_{J}|_{U})\) and
\(\mathcal{G}_{J}|_{U}\) are \emph{quasi-coherent} \(\mathcal{O}_{U}\)-modules,
so is \(\mathcal{G}\) (\cf\ EGA~III, 1.4.17).

Conversely, since one has supposed \(\mathcal{G}_{J}\) \emph{flat} over
\(S_{J}\), if \(\mathcal{G}\) is a \emph{quasi-coherent}
\(\mathcal{O}_{U}\)-module such that
\(\mathcal{G}/\mathcal{J}\mathcal{G}\simeq\mathcal{G}_{J}\) and the morphism
\(\mathcal{J}\otimes_{\mathcal{O}_{S_{J}}}\mathcal{G}_{J}
  \to\mathcal{J}\mathcal{G}\)
is bijective, then \(\mathcal{G}\) is \emph{flat} over \(S\), by the
``fundamental criterion of flatness'' (\cf\ SGA~1~IV,
5.5%
{\renewcommand{\thefootnote}{(78)}\nde{See also \textup{[BAC]},
\S~III.5, th.~1.}}).

Consequently, the set \(\mathcal{E}(U)\) considered here coincides with the
set considered in~4.2, taking for \(f\) the identity morphism of
\(\mathcal{J}\otimes_{\mathcal{O}_{S_{J}}}\mathcal{G}_{J}\), and the
conclusion therefore follows from~4.2.
\end{proof}

{\renewcommand{\thefootnote}{(79)}\nde{One has set out what follows and
added corollary~4.3.1. On the other hand, one recalls that ``pseudo-torsor''
is synonymous with ``formally principal homogeneous''. }}%
One keeps the preceding notations. Let \(Y_{J}\) be a closed subscheme of
\(X_{J}\), defined by a \emph{quasi-coherent} ideal \(\mathcal{I}_{Y_{J}}\).
One supposes \(Y_{J}\) \emph{flat} over \(S_{J}\). Then, applying~4.3 to
\(\mathcal{F}=\mathcal{O}_{X}\) and
\(\mathcal{G}_{J}=\mathcal{O}_{Y_{J}}
  =\mathcal{O}_{X_{J}}/\mathcal{I}_{Y_{J}}\),
one obtains the following corollary.

\begin{corollary}{4.3.1}\label{III.4.3.1}
Let \(S\), \(S_{J}\), \(\mathcal{J}\), \(X\), \(X_{J}\), \(Y_{J}\) and
\(\mathcal{I}_{Y_{J}}\) be as above; one supposes \(Y_{J}\) flat over
\(S_{J}\). Let us denote by \(\mathcal{A}_{J}\) the sheaf of commutative
groups
\[
\mathcal{H}\mathrm{om}_{\mathcal{O}_{X_{J}}}\bigl(\mathcal{I}_{Y_{J}},
  \mathcal{J}\otimes_{\mathcal{O}_{S_{J}}}\mathcal{O}_{Y_{J}}\bigr)
\]
on \(X_{J}\) and \(\mathcal{A}=i_{*}(\mathcal{A}_{J})\), where \(i\) is the
immersion \(X_{J}\hookrightarrow X\).

For every open \(U\) of \(X\), let \(\mathcal{E}(U)\) be the set of closed
subschemes \(Y\) of \(U\), flat over \(S\), such that
\(Y\times_{S}S_{J}=Y_{J}\cap U\). Then \(\mathcal{E}\) is an
\(\mathcal{A}\)-pseudo-torsor.
\end{corollary}

If moreover a \(Y\) exists locally (\ie, if every \(x\in X\) has an open
neighborhood \(U\) such that \(\mathcal{E}(U)\neq\varnothing\)), then
\(\mathcal{E}\) is an \(\mathcal{A}\)-torsor. Now one knows (see for example
EGA~IV\(_{4}\), 16.5.15) that the \(\mathcal{A}\)-torsors on \(X\) are
parametrized by the group
\(H^{1}(X,\mathcal{A})=H^{1}(X_{J},\mathcal{A}_{J})\), and that
\(\mathcal{E}\) possesses a global section (\ie,
\(\mathcal{E}(X)\neq\varnothing\)) if and only if the cohomology class
corresponding to \(\mathcal{E}\) is zero. One therefore obtains the:

\begin{corollary}{4.4}\label{III.4.4}
Let \(S\), \(S_{J}\), \(\mathcal{J}\), \(X\), \(X_{J}\), \(Y_{J}\) and
\(\mathcal{I}_{Y_{J}}\) be as above; one supposes \(Y_{J}\) flat over
\(S_{J}\). Let \(E\) be the set of closed subschemes \(Y\) of \(X\), flat
over \(S\), such that \(Y\times_{S}S_{J}=Y_{J}\).
\begin{enumeratei}
\item The set \(E\) is empty or a principal homogeneous set under the
abelian group
\[
H^{0}(X,\mathcal{A})
  =H^{0}(X_{J},\mathcal{A}_{J})
  =\Hom_{\mathcal{O}_{X_{J}}}\bigl(\mathcal{I}_{Y_{J}},
    \mathcal{J}\otimes_{\mathcal{O}_{S_{J}}}\mathcal{O}_{Y_{J}}\bigr).
\]
\item In order that \(E\) be non-empty, it is necessary and sufficient that
the following two conditions be satisfied:
\begin{enumeratea}
\item There exists locally on \(X\) a solution of the problem.
\item A certain obstruction is zero, which lies in
\[
H^{1}\bigl(X_{J},
  \mathcal{H}\mathrm{om}_{\mathcal{O}_{X_{J}}}\bigl(\mathcal{I}_{Y_{J}},
    \mathcal{J}\otimes_{\mathcal{O}_{S_{J}}}\mathcal{O}_{Y_{J}}\bigr)\bigr).
\]
\end{enumeratea}
\end{enumeratei}
\end{corollary}

\begin{enoncerm}{Complement}{4.4.1}\label{III.4.4.1}
{\renewcommand{\thefootnote}{(80)}\nde{One has added this complement,
useful for proving point~(ii) of proposition~4.8.}}%
Let us keep the notations of~4.4 and suppose that \(E\) contains an element
\(Y\). Let us denote by \(\mathcal{I}_{Y}\) the ideal of \(\mathcal{O}_{X}\)
defining \(Y\), and by \(\mathcal{I}_{Y_{J}}\) its image in
\(\mathcal{O}_{X_{J}}\). Then, as one has seen in the proof of~4.2, one has
a commutative diagram
\[
\xymatrix@C=3.2em@R=2.2em{
\mathcal{J}\otimes_{\mathcal{O}_{S_{J}}}\mathcal{O}_{X_{J}}
  \ar@{->>}[r] \ar[d]
& \mathcal{J}\mathcal{O}_{X} \ar[d] \ar@{-->}[dl] \\
\mathcal{J}\otimes_{\mathcal{O}_{S_{J}}}\mathcal{O}_{Y_{J}}
  \ar[r]^{\sim}
& \mathcal{J}\mathcal{O}_{Y}
}
\]
hence an epimorphism of \(\mathcal{O}_{X}\)-modules
\(\phi\colon\mathcal{J}\mathcal{O}_{X}
  \to\mathcal{J}\otimes_{\mathcal{O}_{S_{J}}}\mathcal{O}_{Y_{J}}\);
let us denote by \(\mathcal{K}\) its kernel. Then, for every element \(Y'\)
of \(E\), the morphism \(\mathcal{O}_{X}\to\mathcal{O}_{Y'}\) factors
through \(\mathcal{O}_{X}/\mathcal{K}\) (which is the amalgamated sum \(B'\)
of the proof of lemma~4.1) and, denoting by \(\mathcal{I}_{Y'}\) the ideal
of \(Y'\) in \(\mathcal{O}_{X}\), one has a commutative diagram:
\[
\xymatrix@C=1.5em@R=1.45em{
& & 0 \ar[d] & 0 \ar[d] & \\
& & \mathcal{I}_{Y'}/\mathcal{K} \ar[r]^{\sim} \ar[d]
  & \mathcal{I}_{Y_{J}} \ar[d] & \\
0 \ar[r]
  & (\mathcal{J}\mathcal{O}_{X})/\mathcal{K} \ar[r] \ar[d]^{\sim}
  & \mathcal{O}_{X}/\mathcal{K} \ar[r] \ar[d]
  & \mathcal{O}_{X_{J}} \ar[r] \ar[d]
  & 0 \\
0 \ar[r]
  & \mathcal{J}\otimes_{\mathcal{O}_{S_{J}}}\mathcal{O}_{Y_{J}}
      \ar[r] \ar[d]
  & \mathcal{O}_{Y'} \ar[r] \ar[d]
  & \mathcal{O}_{Y_{J}} \ar[r] \ar[d]
  & 0 \\
& 0 & 0 & 0 &
}
\]
Thus, replacing \(X\) by the closed subscheme defined by \(\mathcal{K}\), one
reduces to \(\mathcal{K}=0\). Then, giving \(Y'\) is equivalent to giving
the sub-\(\mathcal{O}_{X}\)-module \(\mathcal{I}_{Y'}\) of
\(\mathcal{O}_{X}\), mapping bijectively onto \(\mathcal{I}_{Y_{J}}\) by the
projection \(p\colon\mathcal{O}_{X}\to\mathcal{O}_{X_{J}}\); let us denote
by \(f'\colon\mathcal{I}_{Y_{J}}\isomto\mathcal{I}_{Y'}\)
(\resp \(f\colon\mathcal{I}_{Y_{J}}\isomto\mathcal{I}\)) the inverse
isomorphism.%
% typo?: kept as printed (target of f is I, not I_Y)
Then \(f'-f\) is an element of
\[
\Hom_{\mathcal{O}_{X_{J}}}\bigl(\mathcal{I}_{Y_{J}},
  \mathcal{J}\otimes_{\mathcal{O}_{S_{J}}}\mathcal{O}_{Y_{J}}\bigr)
  =\Hom_{\mathcal{O}_{X_{J}}}\bigl(\mathcal{I}_{Y_{J}},
    \mathcal{J}\mathcal{O}_{Y}\bigr)
\]
which one will denote by \(d(Y',Y)\). (Note that
\(d(Y,Y')=-d(Y',Y)\).)
\end{enoncerm}
